\section{Uso Prático e Ecossistema}

\subsection{Linguagens e Bibliotecas}
O NATS é agnóstico de linguagem, com suporte oficial e comunitário para dezenas de linguagens, incluindo Go, Python, Java, C/C++, C\#, e JavaScript.

\subsection{Ferramentas Disponíveis}

Para o gerenciamento e monitoramento do cluster, o ecossistema disponibiliza ferramentas oficiais:

\begin{itemize}
\item \textbf{NATS CLI}: administração, testes de publicação/assinatura e gerenciamento do JetStream;

\item \textbf{NATS Server}: servidor distribuído em binário único, adequado para ambientes Docker e Kubernetes.
\end{itemize}

\subsection{Aplicações Reais}
Casos reais de adoção mostram o uso do NATS em aplicações distribuídas com requisitos de desempenho e escalabilidade.
\begin{itemize}

\item \textbf{Finanças e Pagamentos:}
uso em mensageria transacional, processamento de eventos e distribuição de dados em tempo real para sistemas financeiros.

\item \textbf{IoT e Telemetria:}
uso em coleta de métricas, comunicação entre dispositivos distribuídos e arquiteturas edge. Empresas como Siemens aparecem entre casos associados ao ecossistema NATS.

\item \textbf{Microsserviços e Alta Concorrência:}
uso como backbone de comunicação entre serviços distribuídos, com aplicações reportadas em empresas como Capital One, Tinder e outras organizações do ecossistema cloud-native.


\begin{itemize}

\item 
\end{itemize}

